\capitulo{7}{Conclusiones y líneas de trabajo futuras}

En este capítulo se recogen las principales conclusiones obtenidas tras el desarrollo del proyecto y se señalan algunas posibles líneas de continuación. El objetivo no es repetir todo el trabajo realizado, sino valorar el resultado alcanzado y los aspectos que podrían ampliarse en el futuro.

\section{Conclusiones}

El desarrollo de este trabajo ha permitido evolucionar una aplicación web ya existente para el diseño de diagramas Entidad-Relación, ampliando sus capacidades y mejorando su coherencia interna. Al no partir de cero, una parte importante del trabajo consistió en comprender el sistema heredado, estabilizar puntos delicados y adaptar su estructura para incorporar nuevas funcionalidades de forma controlada.

Una de las conclusiones principales es que, en una herramienta visual de modelado, no basta con representar correctamente los elementos en pantalla. Cada elemento del diagrama debe estar reflejado también en el modelo interno de la aplicación, ya que esa información se utiliza después para guardar el diagrama, reconstruirlo, validarlo, transformarlo al modelo relacional y generar el script SQL.

A lo largo del proyecto se ha ampliado el soporte del modelo Entidad-Relación incorporando entidades débiles, atributos compuestos y multivaluados, relaciones ternarias con roles y un soporte inicial para jerarquías ISA. Estas ampliaciones no se limitaron a la parte visual, sino que afectaron también a la persistencia, la validación, la transformación relacional y la generación de SQL.

El soporte de jerarquías ISA se planteó de forma limitada y controlada. La aplicación permite representar una generalización con una o varias especializaciones y transforma esta estructura generando una tabla para la entidad general y una tabla para cada especialización, donde la clave heredada actúa como clave primaria y clave foránea. No se incorporaron restricciones más avanzadas como totalidad o parcialidad, especializaciones solapadas o exclusivas, discriminantes o herencia múltiple.

También se revisó la generación de SQL, especialmente el tratamiento de claves foráneas y la forma de expresar el script final. El objetivo fue obtener una salida más clara y coherente con el modelo relacional generado, reduciendo el uso de sentencias \texttt{ALTER TABLE} cuando las dependencias entre tablas pueden resolverse mediante el orden de creación.

Otro resultado relevante es la mejora de la usabilidad del editor. Se reorganizó la interfaz, se revisaron los mensajes de validación y confirmación, se incorporó acceso a información de ayuda dentro de la aplicación y se añadieron operaciones como selección múltiple, ajuste de vista, exportación visual, importación JSON y generación de estructuras de ejemplo.

Durante el desarrollo también se comprobó la importancia de trabajar de forma incremental. Al tratarse de una aplicación heredada y visual, muchos cambios podían afectar a varias partes del sistema al mismo tiempo. Por ello, dividir el trabajo en tareas acotadas y combinar pruebas automáticas con revisión manual permitió avanzar con más seguridad y reducir el riesgo de introducir regresiones.

En conjunto, el resultado obtenido es una herramienta más completa que la aplicación inicial, con mayor soporte para elementos del modelo Entidad-Relación, una salida SQL más cuidada y una interfaz más clara. Su finalidad no es sustituir la explicación teórica ni competir con herramientas profesionales, sino servir como apoyo práctico en un contexto académico.

La revisión de trabajos relacionados permite situar el resultado frente a otras herramientas. En particular, ERDPlus es la alternativa más próxima, ya que también trabaja con diagramas Entidad-Relación, esquemas relacionales y generación SQL. No obstante, UBU E-R App ofrece una propuesta ajustada al alcance docente de este TFG en aspectos como la ausencia de publicidad, la posibilidad de ejecución local, el soporte de relaciones ternarias con roles, la validación orientada al usuario y la integración entre edición visual, modelo interno y generación SQL.

\section{Líneas de trabajo futuras}

A partir del trabajo realizado, pueden plantearse varias líneas de continuación. Algunas están relacionadas con ampliar el soporte del modelo Entidad-Relación, mientras que otras se orientan a reforzar el uso docente de la herramienta o a mejorar su utilidad en escenarios de mayor tamaño.

\begin{itemize}
    \item \textbf{Ampliar el soporte de jerarquías ISA.}  
    La versión actual incorpora un soporte inicial controlado para ISA. Como trabajo futuro se podrían añadir restricciones de totalidad o parcialidad, especializaciones solapadas o exclusivas, discriminantes y distintas estrategias de transformación relacional, como representar toda la jerarquía en una sola tabla o utilizar tablas independientes para la generalización y las especializaciones.

    \item \textbf{Incorporar agregaciones.}  
    Las agregaciones quedan fuera del alcance implementado. Su incorporación requeriría revisar la representación visual, el metamodelo interno, las reglas de validación y la transformación relacional, ya que implican tratar una interrelación como participante de otra interrelación.

    \item \textbf{Mejorar la configuración de tipos de datos y restricciones SQL.}  
    Actualmente, la generación SQL utiliza tipos generales suficientes para el contexto académico. Una línea futura sería permitir configurar tipos concretos, tamaños, obligatoriedad, restricciones adicionales o valores únicos por atributo.

    \item \textbf{Ampliar la gestión de diagramas grandes mediante vistas.}  
    Para modelos con muchas entidades e interrelaciones, podría incorporarse un sistema de vistas que permitiera trabajar con partes del diagrama de forma separada, manteniendo un modelo global común. Esto ayudaría a mejorar la legibilidad en diagramas complejos.

    \item \textbf{Incorporar un banco de ejercicios y corrección automática.}  
    La aplicación podría reforzar su dimensión docente mediante ejercicios integrados. Por ejemplo, ejercicios de diagrama E/R a modelo relacional, ejercicios inversos desde tablas o ejercicios basados en enunciados textuales. Sobre esta base, podrían añadirse mecanismos de corrección automática que comparasen la solución del alumno con una solución esperada y proporcionasen retroalimentación.

    \item \textbf{Permitir rúbricas configurables por el profesor.}  
    Relacionado con la corrección automática, podría incorporarse un sistema de rúbricas en el que el profesor definiera qué aspectos se valoran y con qué peso: entidades, atributos, claves, relaciones, cardinalidades, entidades débiles, transformación relacional o generación SQL.

    \item \textbf{Explorar el apoyo de inteligencia artificial en ejercicios desde enunciado.}  
    En ejercicios basados en texto, podría estudiarse el uso de inteligencia artificial para generar propuestas iniciales, detectar diferencias con una solución de referencia o sugerir posibles errores. Esta línea debería abordarse con prudencia, manteniendo siempre la revisión docente y admitiendo que puede haber varias soluciones válidas.

    \item \textbf{Generar otras representaciones lógicas.}  
    Aunque este TFG se centra en el modelo relacional y SQL, el modelo Entidad-Relación es independiente del modelo lógico final. Por ello, podría estudiarse la generación de otras representaciones, como estructuras orientadas a documentos para bases de datos tipo MongoDB.

    \item \textbf{Relacionar diagramas Entidad-Relación y diagramas de clases UML.}  
    Otra posible ampliación sería generar diagramas de clases UML a partir de diagramas E/R, o realizar el camino inverso en casos controlados. Esta línea permitiría conectar la herramienta con otros modelos utilizados en ingeniería del software.

    \item \textbf{Ampliar la cobertura de pruebas y validaciones.}  
    Futuras ampliaciones funcionales deberían ir acompañadas de nuevas pruebas unitarias, pruebas end-to-end y revisión manual. En particular, sería conveniente aumentar la cobertura de casos complejos de transformación relacional, importación/exportación, reconstrucción visual y generación SQL.
\end{itemize}